\documentclass[12pt,a4paper]{article}
\usepackage[utf8]{inputenc}
\usepackage[T5]{fontenc}
\usepackage[vietnamese]{babel}
\usepackage[margin=1.3cm,top=0.8cm,bottom=1.0cm,includehead,includefoot,headheight=18pt,headsep=6pt,footskip=16pt]{geometry}
\usepackage{amsmath,amssymb,array,tabularx,enumitem,fancyhdr,tikz,xcolor,needspace,multicol}
\usetikzlibrary{calc,angles,quotes,arrows.meta,babel}

\definecolor{ink}{RGB}{30,55,70}
\definecolor{accent}{RGB}{18,82,112}

\setlength{\parindent}{0pt}
\setlength{\parskip}{2pt}

\newcommand{\headerQuote}{%
  \ifcase\numexpr\value{page}-1\relax
    Mathematics is about understanding, not just numbers%
  \or The essence of mathematics lies in its freedom%
  \or Pure mathematics is the poetry of logical ideas%
  \or In mathematics, asking good questions matters most%
  \or Everything around you is built on mathematics%
  \or Mathematics is the queen of the sciences%
  \or Stay with tough problems a little bit longer%
  \or Math gives us hope that every problem has a solution%
  \or Success is the sum of small daily efforts%
  \or Think deeply, check carefully, solve cleanly%
  \or Genius is 1\% inspiration and 99\% perspiration%
  \or Practice makes perfect — keep solving!%
  \or Mathematics is about understanding, not just numbers%
  \or The essence of mathematics lies in its freedom%
  \or Pure mathematics is the poetry of logical ideas%
  \or In mathematics, asking good questions matters most%
  \or Everything around you is built on mathematics%
  \or Mathematics is the queen of the sciences%
  \or Stay with tough problems a little bit longer%
  \or Math gives us hope that every problem has a solution%
  \or Success is the sum of small daily efforts%
  \or Think deeply, check carefully, solve cleanly%
  \or Genius is 1\% inspiration and 99\% perspiration%
  \or Practice makes perfect — keep solving!%
  \or Mathematics is about understanding, not just numbers%
  \or The essence of mathematics lies in its freedom%
  \or Pure mathematics is the poetry of logical ideas%
  \or In mathematics, asking good questions matters most%
  \or Everything around you is built on mathematics%
  \or Mathematics is the queen of the sciences%
  \or Stay with tough problems a little bit longer%
  \or Math gives us hope that every problem has a solution%
  \or Success is the sum of small daily efforts%
  \or Think deeply, check carefully, solve cleanly%
  \or Genius is 1\% inspiration and 99\% perspiration%
  \or Practice makes perfect — keep solving!%
  \or Mathematics is about understanding, not just numbers%
  \or The essence of mathematics lies in its freedom%
  \or Pure mathematics is the poetry of logical ideas%
  \or In mathematics, asking good questions matters most%
  \else Practice makes perfect — keep solving!%
  \fi
}

\pagestyle{fancy}
\fancyhf{}
\lhead{\footnotesize\textit{\color{gray!70!black}\headerQuote}}
\rhead{\small\textbf{\color{ink}Người biên soạn: Dương Gia Bảo}}
\lfoot{}
\rfoot{\small Trang \thepage/10}
\renewcommand{\headrulewidth}{0.4pt}
\renewcommand{\footrulewidth}{0pt}

\fancypagestyle{dapan}{
  \fancyhf{}
  \lhead{\footnotesize\textit{\color{gray!70!black}\headerQuote}}
  \rhead{\small\textbf{\color{ink}Người biên soạn: Dương Gia Bảo}}
  \lfoot{\small\textit{\color{gray!60!black}Đáp án BTVN Hàm số, Hàm số bậc hai \& Vectơ — Toán 10 KNTT}}
  \rfoot{\small Trang \thepage/2}
  \renewcommand{\headrulewidth}{0.4pt}
  \renewcommand{\footrulewidth}{0pt}
}

\newcommand{\chapterbanner}[2]{%
  \par\vspace{3pt}\noindent
  \begin{tikzpicture}
    \node[draw=accent, fill=accent!6, rounded corners=4pt, inner xsep=10pt, inner ysep=6pt, text width=\dimexpr\linewidth-20.8pt\relax, line width=0.8pt] {%
      {\large\bfseries\color{ink}#1}\\[2pt]
      {\small\color{gray!75!black}#2}%
    };
  \end{tikzpicture}\par\nopagebreak\vspace{5pt}
}

\newcommand{\parttitle}[1]{\par\vspace{3pt}\needspace{3\baselineskip}{\normalsize\bfseries\color{accent}#1}\par\nopagebreak\vspace{2pt}}
\newcommand{\four}[4]{\par\nopagebreak\vspace{1.5pt}\noindent\begin{tabularx}{\linewidth}{@{}*{4}{>{\raggedright\arraybackslash}X}@{}}\textbf{A.} #1&\textbf{B.} #2&\textbf{C.} #3&\textbf{D.} #4\end{tabularx}}
\newcommand{\two}[4]{\par\nopagebreak\vspace{1.5pt}\noindent\begin{tabularx}{\linewidth}{@{}*{2}{>{\raggedright\arraybackslash}X}@{}}\textbf{A.} #1&\textbf{B.} #2\\[3.8pt]\textbf{C.} #3&\textbf{D.} #4\end{tabularx}}

\setlist[enumerate,1]{label=\textbf{Câu \arabic*.},leftmargin=1.15cm,labelsep=.15cm,itemsep=2pt,topsep=2pt,parsep=0pt}

\newcommand{\oli}[1]{%
  \par\nopagebreak\vspace{2.5pt}\noindent
  \begin{tikzpicture}
    \draw[step=0.5cm, color=cyan!25!gray!25, line width=0.25pt] (0,0) grid (\linewidth, #1);
    \draw[color=cyan!50!gray!45, line width=0.45pt, rounded corners=2pt] (0,0) rectangle (\linewidth, #1);
    \node[anchor=north west, font=\scriptsize\itshape\color{gray!75}, inner sep=3pt] at (0, #1) {Bài làm:};
  \end{tikzpicture}\par\vspace{3.5pt}
}

\newcolumntype{C}[1]{>{\centering\arraybackslash}m{#1}}


\begin{document}

\chapterbanner{BÀI TẬP VỀ NHÀ: HÀM SỐ, HÀM SỐ BẬC HAI \& CHƯƠNG VECTƠ}{Môn Toán 10 KNTT — Gồm 30 câu Trắc nghiệm ABCD \& 20 Bài Tự luận (có ô li trình bày)}

\parttitle{PHẦN I. CÂU TRẮC NGHIỆM NHIỀU PHƯƠNG ÁN LỰA CHỌN (30 CÂU)}
\textit{Học sinh khoanh tròn vào chữ cái đứng trước phương án trả lời đúng nhất.}

\parttitle{Chủ đề 1. Bài 15: Hàm số (Câu 1 – 8)}

\begin{enumerate}

\needspace{4\baselineskip}\item Tập xác định $D$ của hàm số $y = \dfrac{3x - 1}{\sqrt{5 - x}}$ là
\four{$(-\infty; 5]$.}{$(-\infty; 5)$.}{$(5; +\infty)$.}{$\mathbb{R} \setminus \{5\}$.}

\needspace{4\baselineskip}\item Tập xác định $D$ của hàm số $y = \sqrt{2x + 6} + \dfrac{1}{x - 2}$ là
\two{$(-3; +\infty) \setminus \{2\}$.}{$[-3; +\infty)$.}{$\mathbb{R} \setminus \{2\}$.}{$[-3; +\infty) \setminus \{2\}$.}

\needspace{4\baselineskip}\item Cho hàm số $f(x) = \begin{cases} 2x + 5 & \text{khi } x \ge 1 \\ x^2 - 3 & \text{khi } x < 1 \end{cases}$. Giá trị của biểu thức $f(3) - f(-2)$ bằng
\four{$10$.}{$12$.}{$8$.}{$6$.}

\needspace{4\baselineskip}\item Điểm nào sau đây thuộc đồ thị của hàm số $y = \dfrac{\sqrt{x + 7}}{x - 1}$?
\four{$(1; 0)$.}{$(-7; 1)$.}{$(2; 3)$.}{$(2; -3)$.}

\needspace{4\baselineskip}\item Cho hàm số bậc nhất $y = (2m - 4)x + 3$. Tất cả các giá trị của tham số $m$ để hàm số đồng biến trên $\mathbb{R}$ là
\four{$m > 2$.}{$m \ge 2$.}{$m < 2$.}{$m \ne 2$.}

\needspace{4\baselineskip}\item Cho hàm số $y = f(x)$ có đồ thị là đường gấp khúc đi qua các điểm $A(-4; 2)$, $B(-1; 6)$, $C(2; -1)$ và $D(5; 4)$. Hàm số đồng biến trên khoảng nào sau đây?
\four{$(-1; 2)$.}{$(-4; 2)$.}{$(-1; 5)$.}{$(2; 5)$.}

\needspace{4\baselineskip}\item Tập giá trị của hàm số $y = 2 + \sqrt{16 - x^2}$ là
\four{$[0; 6]$.}{$[2; 6]$.}{$[2; +\infty)$.}{$[-4; 4]$.}

\needspace{4\baselineskip}\item Đồ thị của hàm số $y = \begin{cases} x + 2 & \text{khi } x \ge 0 \\ 3 - 2x & \text{khi } x < 0 \end{cases}$ cắt trục tung tại điểm có tung độ bằng
\four{$3$.}{$0$.}{$2$.}{$-2$.}

\end{enumerate}

\parttitle{Chủ đề 2. Bài 16: Hàm số bậc hai (Câu 9 – 18)}

\begin{enumerate}[resume*]

\needspace{4\baselineskip}\item Tọa độ đỉnh $I$ của parabol $(P): y = 3x^2 - 12x + 8$ là
\four{$I(-2; 44)$.}{$I(2; 4)$.}{$I(4; 8)$.}{$I(2; -4)$.}

\needspace{4\baselineskip}\item Trục đối xứng của parabol $(P): y = -2x^2 - 8x + 5$ là đường thẳng có phương trình
\four{$x = -2$.}{$x = 2$.}{$y = -2$.}{$x = 4$.}

\end{enumerate}

\newpage

\begin{enumerate}[resume*]

\needspace{4\baselineskip}\item Hàm số bậc hai $y = -x^2 + 6x - 4$ đồng biến trên khoảng nào sau đây?
\four{$(3; +\infty)$.}{$(0; +\infty)$.}{$(-\infty; 3)$.}{$(-3; +\infty)$.}

\needspace{4\baselineskip}\item Parabol $(P): y = 2x^2 - 5x + 2$ cắt trục hoành tại hai điểm phân biệt $A$ và $B$. Độ dài đoạn thẳng $AB$ bằng
\four{$2$.}{$\dfrac{3}{2}$.}{$\dfrac{5}{2}$.}{$1$.}

\needspace{4\baselineskip}\item Giá trị nhỏ nhất của hàm số $y = 2x^2 - 4x + 5$ trên đoạn $[0; 3]$ bằng
\four{$3$.}{$5$.}{$11$.}{$1$.}

\needspace{4\baselineskip}\item Giá trị lớn nhất của hàm số $y = -x^2 + 4x + 1$ trên đoạn $[-1; 1]$ bằng
\four{$5$.}{$-4$.}{$1$.}{$4$.}

\needspace{4\baselineskip}\item Biết parabol $(P): y = ax^2 + bx + 4$ đi qua $A(1; 1)$ và có trục đối xứng $x = 2$. Giá trị của tích $a \cdot b$ bằng
\four{$4$.}{$-4$.}{$-3$.}{$3$.}

\needspace{4\baselineskip}\item Cho parabol $(P): y = ax^2 + bx + c$ ($a \ne 0$) quay bề lõm xuống dưới, có đỉnh $I(2; 5)$ và cắt trục tung tại điểm có tung độ bằng $1$. Khẳng định nào sau đây đúng?
\two{$a < 0,\; b < 0,\; c > 0$.}{$a > 0,\; b > 0,\; c > 0$.}{$a < 0,\; b > 0,\; c > 0$.}{$a < 0,\; b > 0,\; c < 0$.}

\needspace{4\baselineskip}\item Số giao điểm của đường thẳng $d: y = 2x + 1$ và parabol $(P): y = x^2 - 2x + 5$ là
\four{$1$.}{$0$.}{$2$.}{$3$.}

\needspace{4\baselineskip}\item Một quả bóng được đá lên theo quỹ đạo parabol $h(t) = -5t^2 + 15t + 2$ (m), với $t \ge 0$ tính bằng giây. Độ cao lớn nhất của quả bóng bằng
\four{$12{,}5\text{ m}$.}{$13{,}25\text{ m}$.}{$15\text{ m}$.}{$11{,}25\text{ m}$.}

\end{enumerate}

\parttitle{Chủ đề 3. Chương IV: Vectơ (Câu 19 – 30)}

\begin{enumerate}[resume*]

\needspace{4\baselineskip}\item Cho ngũ giác đều $ABCDE$. Số vectơ khác vectơ-không có điểm đầu và điểm cuối là các đỉnh của ngũ giác bằng
\four{$10$.}{$25$.}{$15$.}{$20$.}

\needspace{4\baselineskip}\item Cho hình bình hành $ABCD$ tâm $O$. Khẳng định nào sau đây \textbf{sai}?
\two{$\vec{AB} = \vec{DC}$.}{$\vec{OA} + \vec{OC} = \vec{0}$.}{$\vec{AB} + \vec{AD} = \vec{BD}$.}{$\vec{AD} - \vec{AB} = \vec{BD}$.}

\end{enumerate}

\newpage

\begin{enumerate}[resume*]

\needspace{4\baselineskip}\item Cho tam giác $ABC$ vuông tại $A$ có $AB = 5$, $AC = 12$. Độ dài của vectơ $\vec{u} = \vec{AB} + \vec{AC}$ bằng
\four{$17$.}{$13$.}{$7$.}{$\sqrt{119}$.}

\needspace{4\baselineskip}\item Cho hình thoi $ABCD$ cạnh $a$ có góc $\widehat{ABC} = 120^\circ$. Độ dài của vectơ $\vec{v} = \vec{BA} + \vec{BC}$ bằng
\four{$a$.}{$a\sqrt{3}$.}{$a\sqrt{2}$.}{$2a$.}

\needspace{4\baselineskip}\item Cho đoạn thẳng $AB$ có $M$ là trung điểm và $O$ là một điểm bất kì. Khẳng định nào sau đây luôn đúng?
\two{$\vec{OA} + \vec{OB} = \vec{OM}$.}{$\vec{MA} = \vec{MB}$.}{$\vec{OA} - \vec{OB} = 2\vec{OM}$.}{$\vec{OA} + \vec{OB} = 2\vec{OM}$.}

\needspace{4\baselineskip}\item Cho tam giác $ABC$. Gọi $N$ là điểm thuộc cạnh $BC$ sao cho $NC = 3NB$. Biểu diễn vectơ $\vec{AN}$ theo hai vectơ $\vec{AB}$ và $\vec{AC}$ là
\two{$\vec{AN} = \dfrac{1}{4}\vec{AB} + \dfrac{3}{4}\vec{AC}$.}{$\vec{AN} = \dfrac{3}{4}\vec{AB} - \dfrac{1}{4}\vec{AC}$.}{$\vec{AN} = \dfrac{3}{4}\vec{AB} + \dfrac{1}{4}\vec{AC}$.}{$\vec{AN} = \dfrac{2}{3}\vec{AB} + \dfrac{1}{3}\vec{AC}$.}

\needspace{4\baselineskip}\item Trong mặt phẳng tọa độ $Oxy$, cho $\vec{a} = (3; -1)$ và $\vec{b} = (-2; 4)$. Tọa độ của vectơ $\vec{w} = 2\vec{a} - \vec{b}$ là
\four{$(8; -6)$.}{$(4; 2)$.}{$(8; 2)$.}{$(4; -6)$.}

\needspace{4\baselineskip}\item Trong mặt phẳng tọa độ $Oxy$, cho tam giác $ABC$ có $A(1; 4)$, $B(-3; 2)$ và $C(5; -3)$. Tọa độ trọng tâm $G$ của tam giác $ABC$ là
\four{$G(3; 3)$.}{$G(1; 1)$.}{$G(1; -1)$.}{$G(-1; 1)$.}

\needspace{4\baselineskip}\item Trong mặt phẳng tọa độ $Oxy$, cho hai vectơ $\vec{u} = (m; 3)$ và $\vec{v} = (4; -6)$. Hai vectơ $\vec{u}$ và $\vec{v}$ vuông góc với nhau khi và chỉ khi
\four{$m = -2$.}{$m = 2$.}{$m = -\dfrac{9}{2}$.}{$m = \dfrac{9}{2}$.}

\needspace{4\baselineskip}\item Cho tam giác đều $ABC$ cạnh $a$. Giá trị của tích vô hướng $\vec{AB} \cdot \vec{BC}$ bằng
\four{$-\dfrac{a^2}{2}$.}{$\dfrac{a^2}{2}$.}{$-\dfrac{a^2\sqrt{3}}{2}$.}{$\dfrac{a^2\sqrt{3}}{2}$.}

\needspace{4\baselineskip}\item Trong mặt phẳng tọa độ $Oxy$, cho $\vec{x} = (1; -3)$ và $\vec{y} = (2; -1)$. Côsin của góc giữa hai vectơ $\vec{x}$ và $\vec{y}$ bằng
\four{$\dfrac{1}{2}$.}{$-\dfrac{\sqrt{2}}{2}$.}{$\dfrac{\sqrt{2}}{2}$.}{$\dfrac{\sqrt{3}}{2}$.}

\needspace{4\baselineskip}\item Một người kéo một chiếc xe hàng trên mặt đường nằm ngang bằng lực $\vec{F}$ có độ lớn $80\text{ N}$, hướng lực kéo hợp với phương ngang một góc $30^\circ$. Khi xe dịch chuyển được $15\text{ m}$, công của lực $\vec{F}$ thực hiện được là
\four{$600\text{ J}$.}{$600\sqrt{3}\text{ J}$.}{$1200\text{ J}$.}{$400\sqrt{3}\text{ J}$.}

\end{enumerate}

\newpage

\chapterbanner{PHẦN II. BÀI TẬP TỰ LUẬN (20 BÀI)}{(Học sinh trình bày đầy đủ lập luận và lời giải chi tiết vào phần ô li bên dưới mỗi bài)}
\parttitle{Chủ đề 1. Bài 15: Hàm số (Bài 1 – 4)}

\needspace{5\baselineskip}\textbf{Bài 1.} Tìm tập xác định của các hàm số sau:
\begin{enumerate}[label=\alph*), leftmargin=0.6cm, itemsep=1.5pt, topsep=1.5pt]
    \item $y = \dfrac{\sqrt{x + 3}}{x^2 - 4x + 3}$;
    \item $y = \sqrt{4 - x} + \dfrac{x}{\sqrt{2x + 6}}$.
\end{enumerate}
\oli{5.2cm}

\needspace{5\baselineskip}\textbf{Bài 2.} Cho hàm số $f(x) = \begin{cases} 2x^2 - 3x & \text{khi } x \ge 2 \\ \sqrt{6 - x} & \text{khi } -3 \le x < 2 \\ \dfrac{x - 1}{x + 1} & \text{khi } x < -3 \end{cases}$.
\begin{enumerate}[label=\alph*), leftmargin=0.6cm, itemsep=1.5pt, topsep=1.5pt]
    \item Tính giá trị của $f(3)$, $f(-3)$ và $f(-5)$.
    \item Tìm tất cả các giá trị thực của $x$ thuộc $[2; +\infty)$ sao cho $f(x) = 2$.
\end{enumerate}
\oli{5.2cm}

\newpage

\needspace{5\baselineskip}\textbf{Bài 3.} \begin{minipage}[t]{0.56\linewidth}
Cho hàm số $y = f(x)$ xác định trên đoạn $[-3; 5]$ có đồ thị là đường gấp khúc như hình vẽ bên.
\begin{enumerate}[label=\alph*), leftmargin=0.6cm, itemsep=1.5pt, topsep=1.5pt]
    \item Lập bảng biến thiên của hàm số trên $[-3; 5]$ và chỉ rõ các khoảng đồng biến, nghịch biến.
    \item Tìm giá trị lớn nhất, nhỏ nhất của hàm số trên đoạn $[-3; 5]$.
\end{enumerate}
\end{minipage}\hfill
\begin{minipage}[t]{0.41\linewidth}
\vspace{-6pt}\centering
\begin{tikzpicture}[scale=0.42, >=Stealth]
  \draw[step=1, color=gray!22, line width=0.2pt] (-3.5,-2.5) grid (5.5,4.5);
  \draw[->, line width=0.6pt] (-3.8,0) -- (5.8,0) node[below, font=\scriptsize] {$x$};
  \draw[->, line width=0.6pt] (0,-2.8) -- (0,4.8) node[right, font=\scriptsize] {$y$};
  \node[below right, font=\scriptsize] at (0,0) {$O$};
  \draw[dashed, gray!70, line width=0.4pt] (-3,0) -- (-3,2) -- (0,2);
  \draw[dashed, gray!70, line width=0.4pt] (-1,0) -- (-1,-2) -- (0,-2);
  \draw[dashed, gray!70, line width=0.4pt] (2,0) -- (2,4) -- (0,4);
  \draw[dashed, gray!70, line width=0.4pt] (5,0) -- (5,1) -- (0,1);
  \draw[accent, line width=1.25pt] (-3,2) -- (-1,-2) -- (2,4) -- (5,1);
  \foreach \x/\y in {-3/2, -1/-2, 2/4, 5/1} {
    \fill[accent] (\x,\y) circle (2.6pt);
  }
  \node[below, font=\scriptsize] at (-3,0) {$-3$};
  \node[above, font=\scriptsize] at (-1,0) {$-1$};
  \node[below, font=\scriptsize] at (2,0) {$2$};
  \node[below, font=\scriptsize] at (5,0) {$5$};
  \node[right, font=\scriptsize] at (0,-2) {$-2$};
  \node[right, font=\scriptsize] at (0,1) {$1$};
  \node[right, font=\scriptsize] at (0,2) {$2$};
  \node[left, font=\scriptsize] at (0,4) {$4$};
\end{tikzpicture}
\end{minipage}
\oli{5.0cm}

\needspace{5\baselineskip}\textbf{Bài 4.} Tiền cước nước sinh hoạt hàng tháng của một hộ gia đình được tính theo bậc thang: $10\text{ m}^3$ đầu giá $7\,000\text{ đồng/m}^3$; từ trên $10\text{ m}^3$ đến $20\text{ m}^3$ giá $8\,500\text{ đồng/m}^3$; từ trên $20\text{ m}^3$ trở lên giá $10\,000\text{ đồng/m}^3$.
\begin{enumerate}[label=\alph*), leftmargin=0.6cm, itemsep=1.5pt, topsep=1.5pt]
    \item Lập công thức hàm số biểu thị số tiền $T(x)$ (nghìn đồng) theo lượng nước $x\ (\text{m}^3)$ khi $x > 20$.
    \item Tháng 11, gia đình phải trả tổng cộng $255\,000$ đồng tiền nước. Hỏi gia đình đã dùng bao nhiêu $\text{m}^3$ nước?
\end{enumerate}
\oli{5.2cm}

\parttitle{Chủ đề 2. Bài 16: Hàm số bậc hai (Bài 5 – 11)}

\needspace{5\baselineskip}\textbf{Bài 5.} Xác định hàm số bậc hai $y = ax^2 + bx + c$ ($a \ne 0$), biết đồ thị $(P)$ của nó:
\begin{enumerate}[label=\alph*), leftmargin=0.6cm, itemsep=1.5pt, topsep=1.5pt]
    \item Đi qua điểm $A(0; -3)$ và có đỉnh $I(2; -7)$.
    \item Đi qua ba điểm $M(0; 1)$, $N(1; 4)$ và $P(-1; 2)$.
\end{enumerate}
\oli{5.2cm}

\newpage

\needspace{5\baselineskip}\textbf{Bài 6.} Cho hàm số bậc hai $y = -x^2 + 4x - 3$.
\begin{enumerate}[label=\alph*), leftmargin=0.6cm, itemsep=1.5pt, topsep=1.5pt]
    \item Tìm tọa độ đỉnh, trục đối xứng, giao điểm với hai trục tọa độ và lập bảng biến thiên của hàm số.
    \item Tìm giá trị lớn nhất và giá trị nhỏ nhất của hàm số trên đoạn $[-1; 3]$.
\end{enumerate}
\oli{5.2cm}

\needspace{5\baselineskip}\textbf{Bài 7.} \begin{minipage}[t]{0.58\linewidth}
Cho hàm số bậc hai $y = ax^2 + bx + c$ ($a \ne 0$) có đồ thị là parabol $(P)$ như hình vẽ bên.
\begin{enumerate}[label=\alph*), leftmargin=0.6cm, itemsep=1.5pt, topsep=1.5pt]
    \item Xác định công thức của hàm số bậc hai $(P)$.
    \item Tìm tọa độ các giao điểm của $(P)$ và đường thẳng $d: y = 2x - 5$.
\end{enumerate}
\end{minipage}\hfill
\begin{minipage}[t]{0.39\linewidth}
\vspace{-6pt}\centering
\begin{tikzpicture}[scale=0.46, >=Stealth]
  \draw[step=1, color=gray!22, line width=0.2pt] (-0.8,-1.6) grid (4.5,3.8);
  \draw[->, line width=0.6pt] (-1.0,0) -- (4.8,0) node[below, font=\scriptsize] {$x$};
  \draw[->, line width=0.6pt] (0,-1.8) -- (0,4.0) node[right, font=\scriptsize] {$y$};
  \node[below left, font=\scriptsize] at (0,0) {$O$};
  \draw[dashed, gray!70] (2,0) -- (2,-1) -- (0,-1);
  \draw[dashed, gray!70] (2,-1.6) -- (2,3.8);
  \draw[thick, accent, domain=-0.15:4.15, samples=100] plot (\x, {(\x)^2 - 4*(\x) + 3});
  \fill[accent] (2,-1) circle (2.4pt) node[below right, font=\scriptsize] {$I(2;-1)$};
  \fill[ink] (1,0) circle (2pt) node[above right, font=\scriptsize] {$1$};
  \fill[ink] (3,0) circle (2pt) node[above left, font=\scriptsize] {$3$};
  \fill[ink] (0,3) circle (2pt) node[left, font=\scriptsize] {$3$};
  \node[left, font=\scriptsize] at (0,-1) {$-1$};
  \node[above, font=\scriptsize] at (2,0) {$2$};
\end{tikzpicture}
\end{minipage}
\oli{5.0cm}

\needspace{5\baselineskip}\textbf{Bài 8.} Cho đường thẳng $d: y = 2x + m$ và parabol $(P): y = x^2 - 2x + 3$. Tìm tất cả các giá trị của tham số $m$ để $d$ cắt $(P)$ tại hai điểm phân biệt $A, B$ có hoành độ $x_A, x_B$ thỏa mãn $x_A^2 + x_B^2 = 12$.
\oli{5.2cm}

\newpage

\needspace{5\baselineskip}\textbf{Bài 9.} \begin{minipage}[t]{0.58\linewidth}
Một cây cầu vòm có nhịp vòm dạng parabol bề rộng ở mặt sông là $AB = 20\text{ m}$ và chiều cao từ đỉnh vòm xuống mặt sông là $h = 5\text{ m}$ (hình vẽ).
\begin{enumerate}[label=\alph*), leftmargin=0.6cm, itemsep=1.5pt, topsep=1.5pt]
    \item Với hệ trục $Oxy$ như hình vẽ, hãy viết phương trình của vòm parabol.
    \item Một sà lan rộng $8\text{ m}$ đi qua chính giữa gầm cầu. Tìm chiều cao phần nổi tối đa (mét) của sà lan để qua được gầm cầu.
\end{enumerate}
\end{minipage}\hfill
\begin{minipage}[t]{0.39\linewidth}
\vspace{-6pt}\centering
\begin{tikzpicture}[scale=0.24, >=Stealth]
  \fill[cyan!8] (-12,-0.6) rectangle (12,0);
  \draw[->, line width=0.6pt] (-12,0) -- (12.5,0) node[below, font=\scriptsize] {$x$};
  \draw[->, line width=0.6pt] (0,-1.0) -- (0,6.5) node[right, font=\scriptsize] {$y$};
  \draw[thick, accent, domain=-10:10, samples=100] plot (\x, {-0.05*(\x)^2 + 5});
  \draw[line width=0.9pt, fill=gray!25] (-4,0) rectangle (4,4.2);
  \node[font=\tiny] at (0,2.1) {Sà lan};
  \node[below left, font=\scriptsize] at (0,0) {$O$};
  \node[below, font=\scriptsize] at (-10,0) {$A(-10;0)$};
  \node[below, font=\scriptsize] at (10,0) {$B(10;0)$};
  \node[above left, font=\scriptsize, accent] at (0,5) {$5\text{m}$};
  \draw[<->, gray!80!black] (-4,-1.5) -- (4,-1.5) node[midway, below, font=\scriptsize] {$8\text{m}$};
\end{tikzpicture}
\end{minipage}
\oli{4.0cm}

\needspace{5\baselineskip}\textbf{Bài 10.} Một xưởng sản xuất đồ lưu niệm bán mỗi sản phẩm giá $120$ nghìn đồng thì tiêu thụ được $200$ sản phẩm/tháng. Biết chi phí sản xuất mỗi sản phẩm là $50$ nghìn đồng. Khảo sát cho thấy cứ giảm giá bán $5$ nghìn đồng/sản phẩm thì số lượng bán ra tăng thêm $25$ sản phẩm/tháng.
\begin{enumerate}[label=\alph*), leftmargin=0.6cm, itemsep=1.5pt, topsep=1.5pt]
    \item Gọi $x$ là số lần giảm giá $5$ nghìn đồng ($0 \le x < 14$). Lập hàm số biểu diễn lợi nhuận $L(x)$ (nghìn đồng) theo $x$.
    \item Xưởng nên bán giá bao nhiêu nghìn đồng mỗi sản phẩm để lợi nhuận đạt lớn nhất? Tính lợi nhuận lớn nhất đó (triệu đồng).
\end{enumerate}
\oli{4.0cm}

\needspace{5\baselineskip}\textbf{Bài 11.} Từ mặt đất, một vận động viên đá một quả bóng bay lên theo quỹ đạo có độ cao $h(t) = -4{,}9t^2 + 19{,}6t$ (mét), với $t$ (giây) là thời gian chuyển động ($t \ge 0$).
\begin{enumerate}[label=\alph*), leftmargin=0.6cm, itemsep=1.5pt, topsep=1.5pt]
    \item Sau bao lâu kể từ lúc đá thì quả bóng đạt độ cao lớn nhất? Độ cao lớn nhất đó bằng bao nhiêu mét?
    \item Sau bao lâu quả bóng chạm đất trở lại? Khoảng thời gian quả bóng ở độ cao từ $14{,}7\text{ m}$ trở lên là bao nhiêu giây?
\end{enumerate}
\oli{4.0cm}

\newpage

\parttitle{Chủ đề 3. Chương IV: Vectơ (Bài 12 – 20)}

\needspace{5\baselineskip}\textbf{Bài 12.} Cho $6$ điểm bất kì $A, B, C, D, E, F$ và hình bình hành $ABCD$ tâm $O$. Chứng minh rằng:
\begin{enumerate}[label=\alph*), leftmargin=0.6cm, itemsep=1.5pt, topsep=1.5pt]
    \item $\vec{AB} + \vec{CD} + \vec{EF} = \vec{AF} + \vec{CB} + \vec{ED}$.
    \item Với mọi điểm $M$ tùy ý, ta luôn có $\vec{MA} + \vec{MC} = \vec{MB} + \vec{MD} = 2\vec{MO}$.
\end{enumerate}
\oli{5.2cm}

\needspace{5\baselineskip}\textbf{Bài 13.} Cho tam giác $ABC$. Gọi $G$ là trọng tâm tam giác $ABC$, $K$ là điểm trên cạnh $AB$ sao cho $\vec{AK} = \dfrac{1}{3}\vec{AB}$, và $N$ là điểm trên cạnh $AC$ sao cho $\vec{AN} = \dfrac{2}{3}\vec{AC}$.
\begin{enumerate}[label=\alph*), leftmargin=0.6cm, itemsep=1.5pt, topsep=1.5pt]
    \item Phân tích các vectơ $\vec{AG}$ và $\vec{KN}$ theo hai vectơ $\vec{AB}$ và $\vec{AC}$.
    \item Trên tia đối của tia $CB$ lấy điểm $P$ sao cho $CP = \dfrac{1}{3}BC$ (tức $\vec{BP} = \dfrac{4}{3}\vec{BC}$). Chứng minh ba điểm $K, N, P$ thẳng hàng.
\end{enumerate}
\oli{5.2cm}

\needspace{5\baselineskip}\textbf{Bài 14.} Cho hình vuông $ABCD$ cạnh $a$.
\begin{enumerate}[label=\alph*), leftmargin=0.6cm, itemsep=1.5pt, topsep=1.5pt]
    \item Tính độ dài của các vectơ $\vec{u} = \vec{AB} + \vec{AD}$ và $\vec{v} = \vec{AB} - \vec{AC}$.
    \item Tính các tích vô hướng $\vec{AB} \cdot \vec{BD}$ và $(\vec{AB} + \vec{AD}) \cdot (\vec{AB} - \vec{AD})$.
\end{enumerate}
\oli{5.2cm}

\newpage

\needspace{5\baselineskip}\textbf{Bài 15.} Cho tam giác $ABC$ có $AB = 6$, $AC = 8$ và góc $\widehat{BAC} = 60^\circ$.
\begin{enumerate}[label=\alph*), leftmargin=0.6cm, itemsep=1.5pt, topsep=1.5pt]
    \item Tính tích vô hướng $\vec{AB} \cdot \vec{AC}$ và độ dài cạnh $BC$.
    \item Gọi $M$ là trung điểm cạnh $BC$. Tính độ dài đường trung tuyến $AM$ bằng phương pháp vectơ.
\end{enumerate}
\oli{5.5cm}

\needspace{5\baselineskip}\textbf{Bài 16.} Trong mặt phẳng tọa độ $Oxy$, cho ba điểm $A(1; 3)$, $B(-2; 1)$ và $C(4; -1)$.
\begin{enumerate}[label=\alph*), leftmargin=0.6cm, itemsep=1.5pt, topsep=1.5pt]
    \item Tìm tọa độ các vectơ $\vec{AB}$, $\vec{AC}$ và chứng minh ba điểm $A, B, C$ lập thành một tam giác.
    \item Tìm tọa độ điểm $D$ sao cho tứ giác $ABDC$ là hình bình hành.
\end{enumerate}
\oli{5.5cm}

\needspace{5\baselineskip}\textbf{Bài 17.} Trong mặt phẳng tọa độ $Oxy$, cho tam giác $ABC$ có $A(2; 5)$, $B(-1; 2)$ và $C(5; -4)$.
\begin{enumerate}[label=\alph*), leftmargin=0.6cm, itemsep=1.5pt, topsep=1.5pt]
    \item Chứng minh tam giác $ABC$ vuông tại $B$ và tính diện tích tam giác $ABC$.
    \item Tìm tọa độ tâm $I$ và bán kính $R$ của đường tròn ngoại tiếp tam giác $ABC$.
\end{enumerate}
\oli{5.5cm}

\newpage

\needspace{5\baselineskip}\textbf{Bài 18.} Trong mặt phẳng tọa độ $Oxy$, cho hai điểm $A(1; 4)$ và $B(5; -2)$.
\begin{enumerate}[label=\alph*), leftmargin=0.6cm, itemsep=1.5pt, topsep=1.5pt]
    \item Tìm tọa độ điểm $M$ thuộc trục hoành $Ox$ sao cho ba điểm $M, A, B$ thẳng hàng.
    \item Tìm tọa độ điểm $N$ thuộc trục tung $Oy$ sao cho $|\vec{NA} + \vec{NB}|$ đạt giá trị nhỏ nhất. Tính giá trị nhỏ nhất đó.
\end{enumerate}
\oli{5.5cm}

\needspace{5\baselineskip}\textbf{Bài 19.} Trong mặt phẳng tọa độ $Oxy$, cho hai vectơ $\vec{a} = (3; 4)$ và $\vec{b} = (m; 1)$.
\begin{enumerate}[label=\alph*), leftmargin=0.6cm, itemsep=1.5pt, topsep=1.5pt]
    \item Khi $m = 7$, tính tích vô hướng $\vec{a} \cdot \vec{b}$ và số đo góc giữa hai vectơ $\vec{a}$ và $\vec{b}$.
    \item Tìm tất cả các giá trị của tham số $m$ sao cho hai vectơ $\vec{a}$ và $\vec{b}$ hợp với nhau một góc tù.
\end{enumerate}
\oli{5.5cm}

\needspace{5\baselineskip}\textbf{Bài 20.} \begin{minipage}[t]{0.58\linewidth}
Hai lực $\vec{F}_1, \vec{F}_2$ cùng tác dụng vào một vật tại $O$ có độ lớn $|\vec{F}_1| = 30\text{ N}$, $|\vec{F}_2| = 50\text{ N}$, góc giữa hai lực bằng $60^\circ$ (hình vẽ).
\begin{enumerate}[label=\alph*), leftmargin=0.6cm, itemsep=1.5pt, topsep=1.5pt]
    \item Tính độ lớn của hợp lực $\vec{F} = \vec{F}_1 + \vec{F}_2$.
    \item Vật dịch chuyển theo hướng hợp lực $\vec{F}$ một đoạn $12\text{ m}$. Tính công $A$ của hợp lực $\vec{F}$ và công $A_1$ của riêng lực $\vec{F}_1$.
\end{enumerate}
\end{minipage}\hfill
\begin{minipage}[t]{0.39\linewidth}
\vspace{-6pt}\centering
\begin{tikzpicture}[scale=0.45, >=Stealth]
  \coordinate (O) at (0,0);
  \coordinate (F1) at (3,0);
  \coordinate (F2) at (2.5, 4.33);
  \coordinate (F) at (5.5, 4.33);
  \draw[->, very thick, accent] (O) -- (F1) node[below, font=\scriptsize] {$\vec{F}_1$};
  \draw[->, very thick, accent] (O) -- (F2) node[above left, font=\scriptsize] {$\vec{F}_2$};
  \draw[dashed, gray!60] (F1) -- (F) -- (F2);
  \draw[->, line width=1.1pt, ink] (O) -- (F) node[above right, font=\scriptsize] {$\vec{F}$};
  \pic[draw, angle radius=0.45cm, "$60^\circ$", angle eccentricity=1.55, font=\tiny] {angle = F1--O--F2};
  \node[below left, font=\scriptsize] at (O) {$O$};
\end{tikzpicture}
\end{minipage}
\oli{5.2cm}

\newpage
\pagestyle{dapan}
\setcounter{page}{1}

\begin{center}
  {\Large\bfseries\color{ink} BẢNG ĐÁP ÁN TỔNG HỢP (50 CÂU)}\\[2pt]
  {\small\color{gray!75!black} BTVN Toán 10 KNTT — Hàm số, Hàm số bậc hai \& Chương Vectơ}
\end{center}
\vspace{2pt}

\noindent\textbf{PHẦN I. Câu trắc nghiệm nhiều phương án lựa chọn (30 câu)}
\vspace{1.5pt}
\begin{center}
\renewcommand{\arraystretch}{1.15}
\begin{tabular}{|C{0.95cm}|C{0.95cm}|C{0.95cm}|C{0.95cm}|C{0.95cm}|C{0.95cm}|C{0.95cm}|C{0.95cm}|C{0.95cm}|C{0.95cm}|}
\hline
\textbf{1} & \textbf{2} & \textbf{3} & \textbf{4} & \textbf{5} & \textbf{6} & \textbf{7} & \textbf{8} & \textbf{9} & \textbf{10} \\ \hline
\textbf{B} & \textbf{D} & \textbf{A} & \textbf{C} & \textbf{A} & \textbf{D} & \textbf{B} & \textbf{C} & \textbf{D} & \textbf{A} \\ \hline
\hline
\textbf{11} & \textbf{12} & \textbf{13} & \textbf{14} & \textbf{15} & \textbf{16} & \textbf{17} & \textbf{18} & \textbf{19} & \textbf{20} \\ \hline
\textbf{C} & \textbf{B} & \textbf{A} & \textbf{D} & \textbf{B} & \textbf{C} & \textbf{A} & \textbf{B} & \textbf{D} & \textbf{C} \\ \hline
\hline
\textbf{21} & \textbf{22} & \textbf{23} & \textbf{24} & \textbf{25} & \textbf{26} & \textbf{27} & \textbf{28} & \textbf{29} & \textbf{30} \\ \hline
\textbf{B} & \textbf{A} & \textbf{D} & \textbf{C} & \textbf{A} & \textbf{B} & \textbf{D} & \textbf{A} & \textbf{C} & \textbf{B} \\ \hline
\end{tabular}
\end{center}

\vspace{4pt}
\begin{center}
  {\large\bfseries\color{ink} HƯỚNG DẪN GIẢI CHI TIẾT}
\end{center}
\vspace{1pt}

\begin{multicols}{2}
\raggedcolumns
\footnotesize

\textbf{\color{accent}PHẦN I. TRẮC NGHIỆM (30 CÂU)}

\textbf{Câu 1.} ĐK: $5 - x > 0 \Leftrightarrow x < 5 \Rightarrow D = (-\infty; 5)$. Chọn \textbf{B}.

\textbf{Câu 2.} ĐK: $2x + 6 \ge 0 \Leftrightarrow x \ge -3$ và $x \ne 2$. Chọn \textbf{D}.

\textbf{Câu 3.} $f(3) = 2(3) + 5 = 11$; $f(-2) = (-2)^2 - 3 = 1 \Rightarrow 11 - 1 = 10$. Chọn \textbf{A}.

\textbf{Câu 4.} Với $x = 2 \Rightarrow y = \dfrac{\sqrt{2+7}}{2-1} = 3$. Chọn \textbf{C}.

\textbf{Câu 5.} Đồng biến trên $\mathbb{R} \Leftrightarrow 2m - 4 > 0 \Leftrightarrow m > 2$. Chọn \textbf{A}.

\textbf{Câu 6.} Từ $C(2; -1)$ đến $D(5; 4)$, giá trị $y$ tăng từ $-1$ lên $4$ nên đồng biến trên $(2; 5)$. Chọn \textbf{D}.

\textbf{Câu 7.} Vì $0 \le \sqrt{16 - x^2} \le 4$ nên $2 \le y \le 6$. Chọn \textbf{B}.

\textbf{Câu 8.} Giao $Oy$ tại $x = 0 \Rightarrow y = 0 + 2 = 2$. Chọn \textbf{C}.

\textbf{Câu 9.} $x_I = -\dfrac{-12}{6} = 2$, $y_I = 3(4) - 24 + 8 = -4 \Rightarrow I(2; -4)$. Chọn \textbf{D}.

\textbf{Câu 10.} Trục đối xứng $x = -\dfrac{-8}{2(-2)} = -2$. Chọn \textbf{A}.

\textbf{Câu 11.} $a = -1 < 0$, đỉnh $x = 3 \Rightarrow$ đồng biến trên $(-\infty; 3)$. Chọn \textbf{C}.

\textbf{Câu 12.} $2x^2 - 5x + 2 = 0 \Leftrightarrow x = \dfrac{1}{2}$ hoặc $x = 2 \Rightarrow AB = \left|2 - \dfrac{1}{2}\right| = \dfrac{3}{2}$. Chọn \textbf{B}.

\textbf{Câu 13.} Đỉnh $x = 1 \in [0; 3]$. $y(1) = 3, y(0) = 5, y(3) = 11 \Rightarrow \min = 3$. Chọn \textbf{A}.

\textbf{Câu 14.} Đỉnh $x = 2 \notin [-1; 1]$, hàm đồng biến trên $[-1; 1] \Rightarrow \max = y(1) = 4$. Chọn \textbf{D}.

\textbf{Câu 15.} $-\dfrac{b}{2a} = 2 \Rightarrow b = -4a$. Qua $(1; 1) \Rightarrow a - 4a + 4 = 1 \Rightarrow a = 1, b = -4 \Rightarrow ab = -4$. Chọn \textbf{B}.

\textbf{Câu 16.} Bề lõm xuống $\Rightarrow a < 0$; $x_I = -\dfrac{b}{2a} = 2 \Rightarrow b = -4a > 0$; cắt $Oy$ tại $c = 1 > 0$. Chọn \textbf{C}.

\textbf{Câu 17.} $x^2 - 2x + 5 = 2x + 1 \Leftrightarrow (x - 2)^2 = 0 \Leftrightarrow x = 2$ ($1$ giao điểm). Chọn \textbf{A}.

\textbf{Câu 18.} Đỉnh $t = 1{,}5\text{ s} \Rightarrow h_{\max} = -5(2{,}25) + 22{,}5 + 2 = 13{,}25\text{ m}$. Chọn \textbf{B}.

\textbf{Câu 19.} Số vectơ khác $\vec{0}$: $A_5^2 = 5 \times 4 = 20$. Chọn \textbf{D}.

\textbf{Câu 20.} Quy tắc hình bình hành: $\vec{AB} + \vec{AD} = \vec{AC} \ne \vec{BD}$. Chọn \textbf{C}.

\textbf{Câu 21.} $|\vec{AB} + \vec{AC}| = \sqrt{5^2 + 12^2} = 13$. Chọn \textbf{B}.

\textbf{Câu 22.} $\vec{BA} + \vec{BC} = \vec{BD}$. Vì $\widehat{ABC} = 120^\circ \Rightarrow \widehat{BAD} = 60^\circ \Rightarrow \triangle ABD$ đều $\Rightarrow BD = a$. Chọn \textbf{A}.

\textbf{Câu 23.} Tính chất trung điểm: $\vec{OA} + \vec{OB} = 2\vec{OM}$. Chọn \textbf{D}.

\textbf{Câu 24.} $\vec{NC} = -3\vec{NB} \Leftrightarrow \vec{AC} - \vec{AN} = -3(\vec{AB} - \vec{AN}) \Leftrightarrow 4\vec{AN} = 3\vec{AB} + \vec{AC}$. Chọn \textbf{C}.

\textbf{Câu 25.} $\vec{w} = (6; -2) - (-2; 4) = (8; -6)$. Chọn \textbf{A}.

\textbf{Câu 26.} $G\left(\dfrac{1-3+5}{3}; \dfrac{4+2-3}{3}\right) = (1; 1)$. Chọn \textbf{B}.

\textbf{Câu 27.} $\vec{u} \perp \vec{v} \Leftrightarrow 4m + 3(-6) = 0 \Leftrightarrow m = \dfrac{9}{2}$. Chọn \textbf{D}.

\textbf{Câu 28.} $\vec{AB} \cdot \vec{BC} = a \cdot a \cdot \cos 120^\circ = -\dfrac{a^2}{2}$. Chọn \textbf{A}.

\textbf{Câu 29.} $\cos(\vec{x}, \vec{y}) = \dfrac{1(2) + (-3)(-1)}{\sqrt{10}\sqrt{5}} = \dfrac{5}{5\sqrt{2}} = \dfrac{\sqrt{2}}{2}$. Chọn \textbf{C}.

\textbf{Câu 30.} $A = 80 \cdot 15 \cdot \cos 30^\circ = 600\sqrt{3}\text{ J}$. Chọn \textbf{B}.

\columnbreak

\textbf{\color{accent}PHẦN II. TỰ LUẬN (20 BÀI)}

\textbf{Bài 1.} a) $x + 3 \ge 0$ và $x^2 - 4x + 3 \ne 0 \Rightarrow D_a = [-3; +\infty) \setminus \{1; 3\}$.\\
b) $4 - x \ge 0$ và $2x + 6 > 0 \Rightarrow D_b = (-3; 4]$.

\textbf{Bài 2.} a) $f(3) = 2(9) - 9 = 9$; $f(-3) = \sqrt{9} = 3$; $f(-5) = \dfrac{-6}{-4} = \dfrac{3}{2}$.\\
b) Với $x \ge 2$: $2x^2 - 3x = 2 \Leftrightarrow 2x^2 - 3x - 2 = 0 \Leftrightarrow x = 2$ (nhận) hoặc $x = -\dfrac{1}{2}$ (loại). Vậy $x = 2$.

\textbf{Bài 3.} a) Đồng biến trên $(-1; 2)$; nghịch biến trên $(-3; -1)$ và $(2; 5)$.\\
b) $\max_{[-3; 5]} f(x) = 4$ (tại $x=2$); $\min_{[-3; 5]} f(x) = -2$ (tại $x=-1$).

\textbf{Bài 4.} a) $T(x) = 10(7) + 10(8{,}5) + 10(x - 20) = 155 + 10x - 200 = 10x - 45$ (nghìn đồng, $x > 20$).\\
b) $10x - 45 = 255 \Leftrightarrow 10x = 300 \Leftrightarrow x = 30\text{ m}^3$.

\textbf{Bài 5.} a) Đỉnh $I(2; -7) \Rightarrow y = a(x - 2)^2 - 7$. Qua $A(0; -3) \Rightarrow 4a - 7 = -3 \Rightarrow a = 1 \Rightarrow y = x^2 - 4x - 3$.\\
b) Qua $M(0; 1) \Rightarrow c = 1$. Giải hệ từ $N, P$: $a + b = 3$ và $a - b = 1 \Rightarrow a = 2, b = 1 \Rightarrow y = 2x^2 + x + 1$.

\textbf{Bài 6.} a) Đỉnh $I(2; 1)$, trục $x = 2$, cắt $Ox$ tại $(1; 0), (3; 0)$, cắt $Oy$ tại $(0; -3)$. Đồng biến $(-\infty; 2)$, nghịch biến $(2; +\infty)$.\\
b) Trên $[-1; 3]$: $y(2) = 1, y(-1) = -8, y(3) = 0 \Rightarrow \max = 1, \min = -8$.

\textbf{Bài 7.} a) Đỉnh $I(2; -1)$, qua $(0; 3) \Rightarrow y = a(x-2)^2 - 1 \Rightarrow 4a - 1 = 3 \Rightarrow a = 1 \Rightarrow y = x^2 - 4x + 3$.\\
b) $x^2 - 4x + 3 = 2x - 5 \Leftrightarrow x^2 - 6x + 8 = 0 \Rightarrow (2; -1)$ và $(4; 3)$.

\textbf{Bài 8.} $x^2 - 4x + 3 - m = 0$. $\Delta' = m + 1 > 0 \Leftrightarrow m > -1$. Theo Vi-ét: $x_A + x_B = 4, x_A x_B = 3 - m \Rightarrow x_A^2 + x_B^2 = 16 - 2(3 - m) = 10 + 2m = 12 \Rightarrow m = 1$ (thỏa).

\textbf{Bài 9.} a) Đỉnh $(0; 5)$, qua $(10; 0) \Rightarrow y = a x^2 + 5 \Rightarrow 100a + 5 = 0 \Rightarrow a = -\dfrac{1}{20} \Rightarrow y = -\dfrac{1}{20}x^2 + 5$.\\
b) Tại $x = \pm 4 \Rightarrow y = -\dfrac{16}{20} + 5 = 4{,}2\text{ m}$.

\textbf{Bài 10.} a) $L(x) = (70 - 5x)(200 + 25x) = -125x^2 + 750x + 14000$ (nghìn đồng).\\
b) Đỉnh $x = 3 \Rightarrow$ giá bán $120 - 15 = 105$ nghìn đồng; $L_{\max} = L(3) = 15\,125$ nghìn đồng $= 15{,}125$ triệu đồng.

\textbf{Bài 11.} a) Đỉnh $t = \dfrac{19{,}6}{9{,}8} = 2\text{ s} \Rightarrow h_{\max} = 19{,}6\text{ m}$.\\
b) Chạm đất $h(t) = 0 \Rightarrow t = 4\text{ s}$. $h(t) \ge 14{,}7 \Leftrightarrow t^2 - 4t + 3 \le 0 \Leftrightarrow 1 \le t \le 3 \Rightarrow 2\text{ giây}$.

\textbf{Bài 12.} a) $\text{VT} = (\vec{AF} + \vec{FB}) + (\vec{CB} + \vec{BD}) + (\vec{ED} + \vec{DF}) = \vec{AF} + \vec{CB} + \vec{ED} + (\vec{FB} + \vec{BD} + \vec{DF}) = \text{VP}$.\\
b) $O$ trung điểm $AC, BD \Rightarrow \vec{MA} + \vec{MC} = 2\vec{MO}$ và $\vec{MB} + \vec{MD} = 2\vec{MO}$.

\textbf{Bài 13.} a) $\vec{AG} = \dfrac{1}{3}\vec{AB} + \dfrac{1}{3}\vec{AC}$; $\vec{KN} = \vec{AN} - \vec{AK} = -\dfrac{1}{3}\vec{AB} + \dfrac{2}{3}\vec{AC}$.\\
b) $\vec{AP} = \vec{AB} + \dfrac{4}{3}(\vec{AC} - \vec{AB}) = -\dfrac{1}{3}\vec{AB} + \dfrac{4}{3}\vec{AC} \Rightarrow \vec{KP} = -\dfrac{2}{3}\vec{AB} + \dfrac{4}{3}\vec{AC} = 2\vec{KN} \Rightarrow K, N, P$ thẳng hàng.

\textbf{Bài 14.} a) $|\vec{u}| = AC = a\sqrt{2}$; $\vec{v} = \vec{CB} \Rightarrow |\vec{v}| = a$.\\
b) $\vec{AB} \cdot \vec{BD} = a \cdot a\sqrt{2}\cos 135^\circ = -a^2$; $(\vec{AB}+\vec{AD})\cdot(\vec{AB}-\vec{AD}) = AB^2 - AD^2 = 0$.

\textbf{Bài 15.} a) $\vec{AB}\cdot\vec{AC} = 6 \cdot 8 \cdot \cos 60^\circ = 24$; $BC^2 = AB^2 + AC^2 - 2\vec{AB}\cdot\vec{AC} = 52 \Rightarrow BC = 2\sqrt{13}$.\\
b) $AM^2 = \dfrac{1}{4}(\vec{AB}+\vec{AC})^2 = \dfrac{1}{4}(36 + 64 + 48) = 37 \Rightarrow AM = \sqrt{37}$.

\textbf{Bài 16.} a) $\vec{AB} = (-3; -2), \vec{AC} = (3; -4)$. Do $\dfrac{-3}{3} \ne \dfrac{-2}{-4}$ nên $A, B, C$ không thẳng hàng.\\
b) $ABDC$ là hình bình hành $\Leftrightarrow \vec{AB} = \vec{CD} \Rightarrow D = (4 - 3; -1 - 2) = (1; -3)$.

\textbf{Bài 17.} a) $\vec{BA} = (3; 3), \vec{BC} = (6; -6) \Rightarrow \vec{BA}\cdot\vec{BC} = 0 \Rightarrow$ vuông tại $B$. $S = \dfrac{1}{2}(3\sqrt{2})(6\sqrt{2}) = 18$.\\
b) Tâm $I$ trung điểm $AC$: $I\left(\dfrac{7}{2}; \dfrac{1}{2}\right)$; $R = \dfrac{AC}{2} = \dfrac{3\sqrt{10}}{2}$.

\textbf{Bài 18.} a) $M(x; 0)$ thẳng hàng $\Leftrightarrow \dfrac{x - 1}{4} = \dfrac{-4}{-6} \Rightarrow x = \dfrac{11}{3} \Rightarrow M\left(\dfrac{11}{3}; 0\right)$.\\
b) Gọi $I(3; 1)$ trung điểm $AB$: $|\vec{NA} + \vec{NB}| = 2NI$. Nhỏ nhất khi $N$ là hình chiếu lên $Oy \Rightarrow N(0; 1)$, $\min = 6$.

\textbf{Bài 19.} a) $m = 7 \Rightarrow \vec{a}\cdot\vec{b} = 25$; $\cos(\vec{a},\vec{b}) = \dfrac{25}{5 \cdot 5\sqrt{2}} = \dfrac{\sqrt{2}}{2} \Rightarrow 45^\circ$.\\
b) Góc tù $\Leftrightarrow \vec{a}\cdot\vec{b} < 0$ (và không ngược hướng) $\Leftrightarrow 3m + 4 < 0 \Leftrightarrow m < -\dfrac{4}{3}$.

\textbf{Bài 20.} a) $|\vec{F}| = \sqrt{30^2 + 50^2 + 2\cdot 30\cdot 50\cos 60^\circ} = \sqrt{4900} = 70\text{ N}$.\\
b) $A = 70 \cdot 12 = 840\text{ J}$. Ta có $\vec{F}_1\cdot\vec{F} = 30^2 + 30\cdot 50\cos 60^\circ = 1650 \Rightarrow A_1 = \dfrac{12}{70}\cdot 1650 = \dfrac{1980}{7} \approx 282{,}9\text{ J}$.

\end{multicols}
\end{document}
